% 08-privacy.tex: the privacy boundary, stated exactly.
\section{The privacy boundary: what a verifier is allowed to learn}
\label{sec:privacy}

Two things in Polaris use zero knowledge, one thing deliberately does not exist, and the distinction
between them is the most important sentence in this section.

\subsection{Three disclosure levels}

Every verification event is recorded at one of three levels. \code{ZERO\_KNOWLEDGE} proves that a
valid token exists without recording which: the event's token identifier is null, and a check
constraint refuses any row that says otherwise (C2). \code{SELECTIVE} records named attributes only.
\code{FULL} records the identity and is logged and rate-limited. The level is decided by the server
and cannot be upgraded by the client (C6): a relying party cannot turn a zero-knowledge check into a
full one by asking differently, and the Atlas redacts a zero-knowledge event's location server-side.
The default level is zero-knowledge. \sImpl

\subsection{Issuer-side unlinkable verification records}

The first use of zero knowledge is a property of the issuer's database. A default-mode verification
stores no token identifier, so the verification graph, who was verified where and when, cannot be
reconstructed from the database even by someone holding all of it. This is a structural claim, not a
promise: the repository carries a redaction proof that the redacted field is non-derivable under a
stated adversary model, and property tests exercise every read path. The claim is scoped precisely.
It covers zero-knowledge-mode events. Selective and full events do carry a token identifier, because
that is what those levels mean. \sImpl

\subsection{The membership proof}

The second use of zero knowledge is a proof object. A Merkle tree behaves like a cryptographic
fingerprint of an entire set: instead of handing a verifier the whole set, one can hand it a short
proof that a particular item belongs to the set with the committed fingerprint. Polaris commits, per
epoch, a Merkle root over the active-token set, and a holder can produce a Plonky2 proof that its
token is a leaf under that root, bound to the epoch, a context and a nonce, without revealing which
leaf. The verifier sees the proof bytes and four public inputs and cannot recover the leaf; that is
the proof system's property, not an API rule. The setup is transparent, with no ceremony, and the
commitments are hash-based, which is why it is comfortable in a post-quantum threat model. The
default tree depth of fourteen covers the schema's ten-thousand-leaf epoch cap, so the anonymity set
is a full epoch. On the reference machine a proof takes about 35~ms to make and about 13~ms to
check, and the verdict is two-witnessed at the statement level. \sImpl

What the proof does not do is stated with the same care. It proves membership and binding, nothing
else: whether a token is active, permitted for the context and unrevoked is decided when the epoch is
composed, not inside the circuit. The nonce prevents substituting a captured proof into another
context; online, a single-use nonce store refuses a replayed bundle, and offline the verifier keeps no
state, so a relying party that needs single use issues a challenge. Two limitations this section
carried have since been closed and are recorded here as such. The proof now carries a per-verifier
nullifier, \code{Poseidon(secret\,||\,scope\,||\,epoch)}, so a relying party can enforce \emph{one
human, once} within its own scope while another cannot correlate its nullifiers; and the prover runs
on the holder's device, against a published anonymity set, so the issuer is no longer asked which
member is proving. Both rest on the epoch leaf having become a Poseidon commitment the circuit opens:
while it was an opaque digest computed outside, a nullifier beside it would have proved only that the
prover knew some number. \sImpl

Two bounds remain and are not closed. The nullifier does not hide the holder from the \emph{issuer},
which derives every leaf in order to build the tree and could compute any member's nullifier in any
scope; the property is between relying parties. And the circuit over the proof library has had no
external cryptographic audit; the repository's soundness ledger says not to protect real identities
with it as shipped. \sExt

The strongest attack on the proof is not on the proof. It is correlation: matching zero-knowledge
events with selective or full ones across contexts and timestamps to rebuild the graph the proof was
hiding. What answers it is everything around the circuit, the redaction of the graph from every
operator surface, the bounded result sets on every read path, and the threat model's accepted
residual that a sufficiently resourced adversary with many data sources may still deanonymise by
traffic analysis.

\figfile{holder-flow}{Holder, relying party and issuer, with what each sees and what each is
prevented from seeing. The holder presents a file; the relying party decides offline with the
detached verifier or an SDK and learns that a valid credential exists and its value; the issuer
records, in zero-knowledge mode, an event with no token identifier, and keeps no who-verified-whom
row for the online status check. The red edge is the stated limitation: the credential value is a
stable handle across the verifiers it is shown to.}{fig:holder-flow}

\subsection{Bounded, not permanent}

Polaris is a full-credential presentation system, and until recently the consequence was flat: the
login token's subject was the credential's hash, the same value at every relying party, so two of
them could join their user tables exactly and forever without either doing anything wrong. The
identifier they had been handed was a global one.

That is no longer true, and the honest restatement has two halves. The subject and the presentation
handle are now derived under the relying party's own scope, so each is stable where an account needs
it and unrecognisable at the next verifier. Since the subject is the value a relying party
\emph{writes down}, and written-down values are the ones pooled, sold, subpoenaed and breached, this
is the half that matters most in practice. \sImpl

The other half is what did not change. A full-credential presentation still shows the verifier a
stable \code{token\_value}, the issuer's signature and the holder's public key. Two relying parties
who deliberately keep the raw material can still correlate. So the guarantee is about what a verifier
should \emph{store}, not about what it is \emph{shown}, and the detached verifier says which it is
giving: a plain presentation is reported as \code{exposed}, never as unlinkable. Withholding the
material entirely needs the zero-knowledge path, where the handle is the scoped nullifier and the
verifier sees no stable credential at all; that presentation is reported as \code{bounded}.

So cross-verifier correlation is bounded rather than permanent. Blinded presentations, one-time
presentation tokens and anonymous credentials remain unbuilt, and the technologies that would remove
the remaining exposure are discussed in Section~\ref{sec:research}. \emph{Unlinkable by default}
still means the issuer's zero-knowledge-mode records, and still does not mean that a verifier
looking at a full credential learns nothing stable. \sImpl{} for the scoped derivations, with the
residual exposure stated rather than closed.

\subsection{Bounded surfaces}

The remaining privacy mechanisms bound what a surface can return rather than what is stored. Every
map and API aggregate is hard-capped (C8), so the whole population cannot be extracted through the
operator's atlas, and a request for six and a half million cluster bins returns at most five
thousand. The single-entity investigation pages are built to refuse cross-individual link analysis
and are audited: a meta-audit table records who read which audit surface, with what filter, and how
many rows came back, and it is append-only. The authority-structure layer, Athena, is forbidden by
five checks from referencing any natural-person table or column (Section~\ref{sec:cognition}). And
timestamps, the timestamp authority and the exchange gateway are built to learn nothing about what
they touch, which Section~\ref{sec:institutions} describes. \sImpl

\begin{layercard}{Layer card: the privacy boundary}
\qa{What it is}{Three disclosure levels with a CHECK behind the default; server-side enforcement; the redaction proof; the epoch membership proof and its second witness; bounded aggregates; an audited single-entity investigation surface.}
\qa{Why it exists}{A verification log is a surveillance backbone unless it is structurally unable to be one.}
\qa{What it trusts}{The C2 constraint; the proof library's soundness; the redaction proof's adversary model.}
\qa{What it refuses}{A client's claim about its own disclosure level; a token identifier on a zero-knowledge row; an unbounded read; a person node in the authority layer.}
\qa{What it can see}{At zero knowledge: that a valid token was verified at a time in a context. At full: the identity, logged.}
\qa{What it cannot see}{Which token, in zero-knowledge mode, even from the whole database.}
\qa{If it fails}{An identifier on a zero-knowledge row is unstorable; a cap removed fails the build; a person reference in Athena fails five checks. Cross-verifier correlation is not a failure; it is the documented boundary.}
\qa{Checked by}{The C2 CHECK, \code{check\_c6\_atlas\_redacts\_zk\_location}, the redaction property suite, the ZK second-witness differential, \code{check\_athena\_no\_person}.}
\qa{Now or later}{Now. Linkability-resistant presentation is a different system and later.}
\qa{Inside and outside}{Inside: status. Outside: the holders and relying parties who use the credential.}
\end{layercard}

\nextlayer{A credential bounded in what it reveals still has to be used: held, presented, accepted
and logged in with. The next layer is where the credential meets software that is not the issuer's.}
